% Extensión cuántica de la teoría general de la medición.
% Se integra en el mismo capítulo fuente de 01_medida_ambivalencia.

\section{Conditionnement des sections et mise à jour quantique}
\label{sec:medicion-cuantica-genealogica-rev2}
\index{mesure!conditionnement des sections}
\index{instrument quantique!réalisation généalogique}

La théorie générale précédente situe la mesure dans l’état conjoint du système, du lecteur et de l’appareil. La réalisation quantique ajoute une correspondance précise entre trois niveaux : la restriction de l’espace des sections, la projection spectrale et l’instrument complètement positif. L’évolution conjointe relève du couplage ; la publication du résultat relève du lecteur ; la mise à jour relève du conditionnement par la fibre lue.

\subsection{Contexte de mesure \(\mathfrak M_a=(U_a,q_a,\mathcal C_a,m_0)\), application de sortie et fibre conditionnée \(\Omega_{a,y}\)}

Soient \((\Omega_S,\Sigma_S)\) et \((\Omega_M,\Sigma_M)\) les espaces mesurables d’états généalogiques du système et de l’appareil. Un contexte \(a\) est spécifié par
\begin{equation}
 \mathfrak M_a=(U_a,q_a,\mathcal C_a,m_0),
 \label{eq:contexto-medida-rev2}
\end{equation}
où \(m_0\in\Omega_M\) est la préparation de l’appareil,
\[
 U_a:\Omega_S\times\Omega_M
 \longrightarrow\Omega_S\times\Omega_M
\]
est un couplage mesurable, \((Y_a,\Sigma_{Y_a})\) est l’espace mesurable des résultats, \(q_a:\Omega_M\to Y_a\) est un lecteur mesurable de l’aiguille et \(\mathcal C_a\) déclare la compatibilité et l’horizon de prolongement. Lorsqu’un régime réversible est invoqué, on exige en outre que \(U_a\) soit bijectif et bimesurable ; dans la réalisation hilbertienne, la condition correspondante est l’unitarité de \(\widehat U_a\). La sortie est
\begin{equation}
 \operatorname{Out}_a(x)
 =q_a\!\left(\operatorname{pr}_M U_a(x,m_0)\right).
 \label{eq:salida-contexto-medida-rev2}
\end{equation}

Soit \(\Omega_a\in\Sigma_S\) l’ensemble des sections compatibles avec la préparation et le contexte. L’application \(\operatorname{Out}_a\) est mesurable par composition. Pour \(y\in Y_a\) avec \(\{y\}\in\Sigma_{Y_a}\), on définit la fibre mesurable
\begin{equation}
 \Omega_{a,y}
 :=\{x\in\Omega_a:\operatorname{Out}_a(x)=y\}.
 \label{eq:fibra-resultado-medida-rev2}
\end{equation}

\begin{definition}[Mise à jour généalogique]
La mise à jour déterminée par le résultat \(y\) est la restriction
\[
 \Omega_a\longmapsto\Omega_{a,y}.
\]
À profondeur finie, elle conserve exactement les préfixes qui admettent une extension dans \(\Omega_{a,y}\).
\end{definition}

Pour un préfixe \(s\in S_n\), les futurs compatibles après la lecture forment
\[
 \operatorname{Fut}_{a,y}(s)
 =\{x\in\Omega_{a,y}:x|_n=s\}.
\]
La mise à jour modifie la condition de prolongement et conserve le préfixe déjà parcouru ainsi que le registre du couplage.

\begin{proposition}[Probabilité conditionnelle sur la fibre]
\label{prop:probabilidad-condicionada-medida-rev2}
Soit \(\mu_a\) une probabilité sur l’espace mesurable restreint \((\Omega_a,\Sigma_S|_{\Omega_a})\), soit \(B\in\Sigma_S|_{\Omega_a}\) et supposons \(\mu_a(\Omega_{a,y})>0\). Alors
\begin{equation}
 \mu_{a,y}(B)
 =\frac{\mu_a(B\cap\Omega_{a,y})}
        {\mu_a(\Omega_{a,y})}
 \label{eq:bayes-genealogico-rev2}
\end{equation}
définit une probabilité concentrée sur la fibre du résultat : \(\mu_{a,y}(\Omega_{a,y})=1\).
\end{proposition}

\begin{proof}
L’application est positive ou nulle et dénombrablement additive par les propriétés de \(\mu_a\). La normalisation s’obtient en substituant \(B=\Omega_{a,y}\) dans \eqref{eq:bayes-genealogico-rev2}.
\end{proof}

\subsection{Réalisation spectrale de l’appareil}

Soit \(\{|r\rangle\}\) une base orthonormale adaptée à l’observable du système, soit \(|m_0\rangle\) la préparation de l’appareil et soit \(\{|m_r\rangle\}\) une famille orthonormale d’états d’aiguille. La représentation unitaire du couplage satisfait
\begin{equation}
 \widehat U_a(|r\rangle\otimes|m_0\rangle)
 =|r\rangle\otimes|m_r\rangle.
 \label{eq:premedicion-espectral-rev2}
\end{equation}
Pour \(|\psi\rangle=\sum_r c_r|r\rangle\), la linéarité produit
\[
 \widehat U_a(|\psi\rangle\otimes|m_0\rangle)
 =\sum_r c_r|r\rangle\otimes|m_r\rangle.
\]
Le couplage construit ainsi une corrélation entre le système et l’aiguille. La lecture sélective du résultat \(r\) est représentée, pour un état densité avec \(\operatorname{tr}(P_r\rho)>0\), par la règle de Lüders~\cite{Lueders1950},
\begin{equation}
 \rho\longmapsto
 \rho_r=\frac{P_r\rho P_r}{\operatorname{tr}(P_r\rho)}.
 \label{eq:luders-genealogico-rev2}
\end{equation}

\begin{theorem}[Correspondance entre fibres et projections]
\label{thm:fibra-proyeccion-medida-rev2}
Supposons que chaque fibre \(\Omega_{a,y}\) soit une réunion finie de cylindres disjoints de même profondeur \(n\), que la réalisation \(\mathcal Q\) envoie ces cylindres sur des sous-espaces orthogonaux et que \(P_y\) soit la somme de leurs projecteurs cylindriques. Soit \(\mu_a(\Omega_{a,y})>0\) et soit \(\Psi_n(\mu_a)\) l’état construit par \eqref{eq:estado-amplitudes-medida-rev2} à partir de cette même mesure. Alors
\[
 \mu_a(\Omega_{a,y})=\|P_y\Psi_n(\mu_a)\|^2,
\]
et la restriction à \(\Omega_{a,y}\) est représentée par la projection normalisée \(P_y\Psi_n(\mu_a)/\|P_y\Psi_n(\mu_a)\|\). Dans la carte des matrices densité, si \(\operatorname{tr}(\rho P_y)>0\), soit \(\mathcal A_n\) l’algèbre des événements cylindriques de profondeur \(n\) et supposons que \(B\mapsto P_B\) et \(P_y\) appartiennent à la même PVM et que
\begin{equation}
 \mu_a(B)=\operatorname{tr}(\rho P_B)
 \qquad\text{pour tout }B\in\mathcal A_n.
 \label{eq:equivalencia-medida-pvm-rev2}
\end{equation}
Alors la même probabilité conditionnelle est représentée par \eqref{eq:luders-genealogico-rev2} : pour tout \(B\in\mathcal A_n\),
\[
 \mu_{a,y}(B)=\operatorname{tr}(\rho_yP_B).
\]
\end{theorem}

\begin{proof}
La fibre est une réunion disjointe de cylindres et \(P_y\) est la somme orthogonale de leurs projecteurs. \cref{thm:born-cilindrico-rev2} identifie la norme quadratique de chaque composante à sa mesure ; l’additivité donne l’égalité. Après projection, les coefficients de chaque cylindre sont divisés par \(\sqrt{\mu_a(\Omega_{a,y})}\) ; leurs carrés sont donc les probabilités conditionnelles de \eqref{eq:bayes-genealogico-rev2}. Dans la carte des densités, l’appartenance à la même PVM donne \(P_BP_y=P_{B\cap\Omega_{a,y}}\). Par cyclicité de la trace et par \eqref{eq:equivalencia-medida-pvm-rev2},
\[
 \operatorname{tr}(\rho_yP_B)
 =\frac{\operatorname{tr}(\rho P_{B\cap\Omega_{a,y}})}
        {\operatorname{tr}(\rho P_y)}
 =\frac{\mu_a(B\cap\Omega_{a,y})}{\mu_a(\Omega_{a,y})}.
\]
\end{proof}

\subsection{Instruments induits par le couplage}

Pour une base d’aiguilles raffinée \(\{|m_{y,\alpha}\rangle\}_{y,\alpha}\), les opérateurs
\begin{equation}
 M_{y,\alpha}=\langle m_{y,\alpha}|\widehat U_a|m_0\rangle,
 \qquad
 \sum_{y,\alpha} M_{y,\alpha}^\dagger M_{y,\alpha}=I
 \label{eq:kraus-acoplamiento-genealogico-rev2}
\end{equation}
forment une représentation de Kraus~\cite{Kraus1983}. L’effet \(E_y=\sum_\alpha M_{y,\alpha}^\dagger M_{y,\alpha}\) détermine la probabilité et, lorsque \(p(y\mid a,\rho)>0\), l’état ultérieur :
\begin{equation}
 p(y\mid a,\rho)=\operatorname{tr}(\rho E_y),
 \qquad
 \rho_{a,y}
 =\frac{\sum_\alpha M_{y,\alpha}\rho M_{y,\alpha}^\dagger}
 {p(y\mid a,\rho)}.
 \label{eq:instrumento-kraus-genealogico-rev2}
\end{equation}
La transformation non sélective est
\[
 \mathcal E_a(\rho)=
 \sum_{y,\alpha}M_{y,\alpha}\rho M_{y,\alpha}^\dagger.
\]

\begin{proposition}[Couplage unitaire et canal réduit]
\label{prop:acoplamiento-canal-reducido-rev2}
La dynamique conjointe \(\widehat U_a\) et la préparation \(|m_0\rangle\) induisent sur le système un canal complètement positif préservant la trace. La décomposition selon les aiguilles induit l’instrument \(\mathcal I_{a,y}(\rho)=
\sum_\alpha M_{y,\alpha}\rho M_{y,\alpha}^\dagger\). Réciproquement, tout canal complètement positif admet une dilatation par un espace auxiliaire et une isométrie.
\end{proposition}

\begin{proof}
L’identité de complétude dans \eqref{eq:kraus-acoplamiento-genealogico-rev2} préserve la trace, et la forme de Kraus démontre la positivité complète. L’énoncé réciproque est le théorème de dilatation de Stinespring~\cite{Stinespring1955}.
\end{proof}

L’espace auxiliaire de la dilatation peut réaliser linéairement la mémoire et le registre transportés par TPK dès lors qu’un plongement entrelaçant cette mémoire avec l’espace auxiliaire est déclaré. Stinespring garantit la dilatation du canal ; il n’identifie pas à lui seul les deux objets. La règle de Lüders est le cas projectif fin \(M_{y,1}=P_y\) ; les lecteurs dégénérés et les appareils non idéaux sont décrits par des familles \(M_{y,\alpha}\).

\subsection{Observateur interne et complétude de la lecture}

Un observateur réalisé est un sous-système capable de préparer un contexte, de se coupler à une frontière, de conserver un registre discernable et de l’utiliser comme condition de prolongement. La suite
\begin{equation}
 (x,m_0)
 \xrightarrow{\ U_a\ }(x',m')
 \xrightarrow{\ q_a\ }y
 \xrightarrow{\ \mathrm{cond}\ }\Omega_{a,y}
 \label{eq:secuencia-observador-rev2}
\end{equation}
situe le cadre de lecture et la mémoire de l’observateur dans l’état élargi.

Le principe d’Ambivalence complète cette description. Les lectures aréale et volumétrique produisent des quotients \(q_A,q_V\), tandis que la mémoire \(T\) conserve le transport entre elles. Lorsque
\[
 \Xi:x\longmapsto(q_A(x),q_V(x),T(x))
\]
est injective modulo le groupe de jauge déclaré, le couple d’observables et la mémoire reconstruisent la classe de l’état complet.

\begin{theorem}[Décomposition généalogique de la mesure]
\label{thm:descomposicion-medicion-rev2}
Sous les hypothèses de \cref{thm:fibra-proyeccion-medida-rev2,prop:acoplamiento-canal-reducido-rev2}, la transition entre évolution et résultat se décompose, dans le régime où \(U_a\) est bijectif et bimesurable ou \(\widehat U_a\) est unitaire, en :
\[
 \text{couplage réversible}
 \longrightarrow
 \text{lecture du pointeur}
 \longrightarrow
 \text{conditionnement des sections}.
\]
Cette décomposition comporte deux régimes :
\begin{enumerate}
\item Dans une mesure cylindrique projective, \(M_{y,1}=P_y\), et la restriction de fibre, la projection normalisée et la mise à jour de Lüders coïncident :
\[
 \Omega_a\longmapsto\Omega_{a,y},
 \qquad
 \Psi_n\longmapsto\frac{P_y\Psi_n}{\|P_y\Psi_n\|},
 \qquad
 \rho\longmapsto
 \frac{P_y\rho P_y}{\operatorname{tr}(\rho P_y)}.
\]
\item Pour un instrument général, l’état conditionné est
\[
 \rho\longmapsto
 \frac{\sum_\alpha M_{y,\alpha}\rho M_{y,\alpha}^\dagger}
 {\operatorname{tr}(\rho E_y)}.
\]
Dans une dilatation, cette mise à jour provient d’une projection de l’aiguille sur le système composé avec l’appareil. Sa correspondance avec \(\Omega_{a,y}\) exige un morphisme déclaré qui entrelace la fibre du lecteur, l’effet \(E_y\) et l’instrument \(\mathcal I_{a,y}\).
\end{enumerate}
\end{theorem}

\begin{proof}
La première flèche réalise le couplage ; la troisième restreint l’espace des sections à la fibre définie dans \eqref{eq:fibra-resultado-medida-rev2}. \cref{thm:fibra-proyeccion-medida-rev2} démontre la coïncidence des trois formes dans le régime projectif. Dans le régime général, \eqref{eq:instrumento-kraus-genealogico-rev2} démontre la mise à jour de l’instrument et Stinespring fournit sa réalisation dilatée. Le morphisme d’entrelacement supplémentaire est précisément la donnée qui transporte cette mise à jour en retour vers l’espace des sections.
\end{proof}

\begin{statusbox}[title={Résultat et localisation}]
Dans le modèle déclaré, la mise à jour quantique est la réalisation spectrale de la restriction des futurs compatibles. La réversibilité se décide par la bimesurabilité de \(U_a\) ou l’unitarité de \(\widehat U_a\) ; la perte d’information, dans les fibres du lecteur ou dans un effacement ultérieur ; et l’état conditionné, dans \(\Omega_{a,y}\). Le coût thermodynamique de réinitialisation et le coût de Landauer nul de l’évolution réversible demeurent dans le chapitre consacré à l’information et à l’entropie.
\end{statusbox}
